\begin{enumerate}
\item 
	\begin{enumerate}
		\item Les vecteurs $\overrightarrow{\text{AB}}$ et $\overrightarrow{\text{AC}}$ ont pour coordonnées :
\[
\begin{array}{rcl}
\overrightarrow{\text{AB}} && \begin{pmatrix} -1 - 1 \\ 3 - 2 \\ 1 - 3 \end{pmatrix} = \begin{pmatrix} -2 \\ 1 \\ -2 \end{pmatrix} \\
\overrightarrow{\text{AC}} && \begin{pmatrix} 2 - 1 \\ 1 - 2 \\ 6 - 3 \end{pmatrix} = \begin{pmatrix} 1 \\ -1 \\ 3 \end{pmatrix}
\end{array}
\]
Leurs coordonnées ne sont pas proportionnelles. Les vecteurs ne sont pas colinéaires, les points A, B et C ne sont pas alignés et définissent donc un plan.

		\item Calculons les produits scalaires :
\[
\begin{array}{rcl}
\vect{n} \cdot \vect{\text{AB}} &=& 1 \times (-2) + 4 \times 1 + 1 \times (-2) = -2 + 4 - 2 = 0 \\
\vect{n} \cdot \vect{\text{AC}} &=& 1 \times 1 + 4 \times (-1) + 1 \times 3 = 1 - 4 + 3 = 0
\end{array}
\]
Le vecteur $\vect{n}$ est orthogonal à deux vecteurs non colinéaires du plan (ABC), il est donc normal à ce plan.

\item Une équation du plan (ABC) est de la forme $x + 4 y + z + d = 0$. Le point A$(1~;~2~;~3)$ appartient au plan, donc :
\[
\begin{array}{rcl}
1 + 4\times 2+ 3 + d &=& 0 \iff 12 + d = 0 \iff d = -12
\end{array}
\]
Une équation cartésienne est : $x + 4y + z - 12 = 0$.
\end{enumerate}
\item 
\begin{enumerate}
\item La droite $(d)$ est dirigée par le vecteur normal $\vect{n}(1~;~4~;~1)$ et passe par D$(3~;~-2~;~-1)$. Une représentation paramétrique est :
\[
\left\{
\begin{array}{rcr}
x &=& 3 + \phantom{4}t \\
y &=& -2 + 4t \\
z &=& -1 + \phantom{4}t
\end{array}
\right. , t \in \mathbb{R}
\]

\item Le point H appartient à $(d)$, il existe donc $t \in \mathbb{R}$ tel que $\left\{ \begin{matrix} x_{\text{H}}&=&3+\phantom{4}t \\y_{\text{H}}&=& -2+4t \\ z_{\text{H}}&=& -1+\phantom{4}t \end{matrix} \right.$. De plus H appartient au plan (ABC), ses coordonnées vérifient donc l'équation du plan (ABC): $x_{\text{H}}+4y_{\text{H}}+z_{\text{H}}- 12 = 0$. On en déduit que $t$ est solution de :
\[
\begin{array}{rcl}
(3 + t) + 4(-2 + 4t) + (-1 + t) - 12 &=& 0 \\
3 + t - 8 + 16t - 1 + t - 12 &=& 0 \\
18t - 18 &=& 0 \\
t &=& 1
\end{array}
\]
Les coordonnées de H sont donc $x_{\text{H}} = 3 + 1 = 4$, $y_{\text{H}} = -2 + 4(1) = 2$, $z_{\text{H}} = -1 + 1 = 0$. Soit H$(4~;~2~;~0)$.

\item La distance de D au plan (ABC) correspond à la longueur DH :
\[
\begin{array}{rcl}
\text{DH} &=& \sqrt{(4 - 3)^2 + (2 - (-2))^2 + (0 - (-1))^2} \\
&=& \sqrt{1^2 + 4^2 + 1^2} \\
&=& \sqrt{18} = 3\sqrt{2}
\end{array}
\]
\end{enumerate}

\item 
\begin{enumerate}
\item On sait que $\vect{\text{AB}} \cdot \vect{\text{AC}} = \text{AB} \times \text{AC} \times \cos(\widehat{\text{BAC}})$.

\[
\begin{array}{rcl}
\overrightarrow{\text{AB}} \cdot \overrightarrow{\text{AC}} &=& -2(1) + 1(-1) - 2(3) = -9 \\
\text{AB} &=& \sqrt{(-2)^2 + 1^2 + (-2)^2} = \sqrt{9} = 3 \\
\text{AC} &=& \sqrt{1^2 + (-1)^2 + 3^2} = \sqrt{11} \\
\end{array}
\]
D'où :
\[
\begin{array}{rcl}
-9 &=& 3 \times \sqrt{11} \times \cos(\widehat{\text{BAC}}) \\
\cos(\widehat{\text{BAC}}) &=& \dfrac{-9}{3\sqrt{11}} = \dfrac{-3}{\sqrt{11}} = \dfrac{-3\sqrt{11}}{11}
\end{array}
\]

\item Sachant que $\cos^2(\widehat{\text{BAC}}) + \sin^2(\widehat{\text{BAC}}) = 1$ :
\[
\begin{array}{rcl}
\sin^2(\widehat{\text{BAC}}) &=& 1 - \left(-\frac{3}{\sqrt{11}}\right)^2 \\
&=& 1 - \frac{9}{11} = \dfrac{2}{11}
\end{array}
\]
L'angle $\widehat{\text{BAC}}$ étant dans $[0~;~\pi]$, son sinus est positif : $\sin(\widehat{\text{BAC}}) = \dfrac{\sqrt{2}}{\sqrt{11}}$.

\item L'aire du triangle ABC est donnée par :
\[
\begin{array}{rcl}
\mathcal{A} &=& \frac{1}{2} \times \text{AB} \times \text{AC} \times \sin(\widehat{\text{BAC}}) \\
&=& \frac{1}{2} \times 3 \times \sqrt{11} \times \dfrac{\sqrt{2}}{\sqrt{11}} \\
&=& \dfrac{3\sqrt{2}}{2}
\end{array}
\]
\end{enumerate}

\item Le volume du tétraèdre ABCD est :
\[
\begin{array}{rcl}
\mathcal{V} &=& \dfrac{1}{3} \times \mathcal{A}_{\text{ABC}} \times \text{DH} \\
&=& \frac{1}{3} \times \dfrac{3\sqrt{2}}{2} \times 3\sqrt{2} \\
&=& \frac{1}{3} \times \dfrac{18}{2} \\
&=& 3
\end{array}
\]
\end{enumerate}
